In 2021, in response to the 2020 update of the European Strategy for Particle Physics, the CERN Council initiated the Future Circular Collider (FCC) Feasibility Study. 

This report summarises an immense amount of work carried out by the international FCC collaboration over several years. It covers, inter alia, physics objectives and potential, geology, civil engineering, technical infrastructure, territorial implementation, environmental aspects, R\&D needs for the accelerators and detectors, socio-economic benefits and cost. It constitutes important input for the ongoing update of the European Strategy for Particle Physics.

The Feasibility Study required engagement with a broad range of stakeholders. In particular, throughout the Study, CERN has been accompanied by its two Host States, France and Switzerland, and has been working with entities at local, regional and national level. I am very grateful to the Host State authorities and teams for their invaluable help. Furthermore, significant sections of the Study were supported by the European Union under the Horizon 2020 and Horizon Europe framework programmes. The Study also greatly benefited from contributions from accelerator laboratories and universities from across Europe, such as the Swiss Accelerator Research and Technology (CHART) initiative, and from the Americas, Asia, Africa and Australia. 

The proposed FCC integrated programme consists of two possible stages: an electron–positron collider serving as a Higgs-boson, electroweak and top-quark factory running at different centre-of-mass energies, followed at a later stage by a proton–proton collider operating at an unprecedented collision energy of around 100 TeV. The complementary physics programmes of each stage match the physics priorities expressed in the 2020 update of the European Strategy for Particle Physics. 

A major achievement of the Feasibility Study is the choice of placement of the collider ring and the entire infrastructure, including the surface sites and the access shafts, which was developed and optimised over several years following the principle “avoid, reduce, compensate”. Sustainability studies have assessed energy efficiency, land use, water and resource management, and socio-economic impact, ensuring that the FCC is designed in accordance with the latest environmental and societal standards. 

I would like to thank all contributors to this report for their hard work and commitment, which allowed the outstanding results presented here to be achieved. 


\begin{flushright}
\textbf{Fabiola Gianotti}\\
CERN, Director-General
\end{flushright}